\documentclass[12pt,a4paper]{article}

% =========================================================
% PACOTES
% =========================================================
\usepackage[utf8]{inputenc}
\usepackage[T1]{fontenc}
\usepackage[brazil]{babel}
\usepackage{geometry}
\usepackage{graphicx}
\usepackage{booktabs}
\usepackage{tabularx}
\usepackage{longtable}
\usepackage{array}
\usepackage{enumitem}
\usepackage{xcolor}
\usepackage{hyperref}
\usepackage{fancyhdr}
\usepackage{setspace}
\usepackage{float}

\geometry{
    top=2.5cm,
    bottom=2.5cm,
    left=2.5cm,
    right=2.5cm
}

\setstretch{1.15}
\setlength{\headheight}{15pt}

\definecolor{titulo}{RGB}{45,70,90}

\hypersetup{
    colorlinks=true,
    linkcolor=titulo,
    urlcolor=titulo
}

\pagestyle{fancy}
\fancyhf{}
\fancyhead[L]{\small Ecos do Cretáceo: Projeto Retorno}
\fancyhead[R]{\small Game Design Document}
\fancyfoot[C]{\thepage}

\setlist[itemize]{leftmargin=1.2cm}
\setlist[enumerate]{leftmargin=1.2cm}

\newcommand{\jogo}{\textbf{Ecos do Cretáceo: Projeto Retorno}}

% =========================================================
% DOCUMENTO
% =========================================================

\begin{document}

% =========================================================
% 1. PÁGINA DE TÍTULO
% =========================================================

\begin{titlepage}
    \centering

    \vspace*{2cm}

    {\Huge\bfseries ECOS DO CRETÁCEO\par}
    \vspace{0.3cm}
    {\LARGE\bfseries PROJETO RETORNO\par}

    \vspace{0.8cm}

    {\Large Game Design Document -- GDD\par}

    \vspace{1.5cm}

    {\large
    Um jogo 3D de aventura, exploração, furtividade e sobrevivência
    sobre um humano sequestrado por uma civilização de dinossauros
    tecnologicamente avançada.
    \par}

    \vfill

    \begin{flushleft}
        \textbf{Disciplina:} Desenvolvimento de Games\\
        \textbf{Professor:} Saulo Pereira Ribeiro\\
        \textbf{Curso:} Ciência da Computação\\
        \textbf{Instituição:} Universidade de Vila Velha\\[0.5cm]

        \textbf{Integrantes:}\\
        Isabel Sterim Saade\\
        Maria Carla dos Santos Bellote\\
        Matheus Pastore Morgan\\
       
    \end{flushleft}

    \vfill

    {\large Revisão de 21 de outubro de 2026\par}

    \thispagestyle{empty}
\end{titlepage}

\addcontentsline{toc}{section}{1. Página de Título}

\tableofcontents
\clearpage

% Faz a próxima section começar no número 2.
\setcounter{section}{1}

% =========================================================
% 2. VISÃO GERAL
% =========================================================

\section{Visão geral do Jogo}

\subsection*{Estado do documento e da implementação}

Esta revisão, de 21 de outubro de 2026, descreve o \textbf{Nível 1
já jogável} e o \textbf{planejamento da campanha com cinco níveis}.
Os níveis 2 a 5, seus personagens, diálogos, combates e eventos
continuam sendo propostas de design.

A versão atual contém Menu Principal, Nível 1 com progressão completa,
Menu de Pausa e tela de conclusão. Ainda não existe transição jogável
para o Nível 2. As cenas da seção narrativa orientam o desenvolvimento
e não representam conteúdo já encenado no jogo.

\begin{center}
\begin{tabularx}{\textwidth}{|p{3cm}|X|}
\hline
\textbf{Situação} & \textbf{Conteúdo}\\
\hline
Implementado & Movimentação, corrida, salto, agachamento, interação F,
arremesso, bateria, cartão aleatório, três terminais, portas, HUD,
menus, pet adotável, partículas e áudio disponível.\\
\hline
Preparado & Componentes de passos e cantos do pássaro, ainda sem
AudioClips adequados atribuídos.\\
\hline
Planejado & Níveis 2--5, narrativa encenada, inimigos, furtividade,
combate, vida, checkpoints, salvamento, diálogos e NavMesh.\\
\hline
\end{tabularx}
\end{center}



\subsection{Público alvo}

O público-alvo de \jogo{} é composto principalmente por adolescentes
e jovens adultos, aproximadamente entre 14 e 30 anos, interessados em
jogos de aventura, ficção científica, mistério, exploração e
sobrevivência.

O jogo não pretende utilizar violência gráfica extrema. O conflito
será representado principalmente por perseguições, furtividade,
combate leve e situações de tensão.

A classificação indicativa pretendida é de aproximadamente 12 ou
14 anos, dependendo da intensidade final das cenas de combate.

\subsection{Conceito do Jogo}

\jogo{} é um jogo 3D em primeira pessoa desenvolvido na Unity.

O jogador controla Alex Moraes, um humano que desperta dentro de uma
câmara de contenção em uma nave desconhecida sem saber onde está,
quem o sequestrou ou quanto tempo permaneceu desacordado.

Inicialmente, tudo indica que Alex foi abduzido por uma espécie
alienígena.

Ao escapar da nave, entretanto, ele descobre que está em um planeta
distante e que seus sequestradores são, na verdade, descendentes de
dinossauros que desenvolveram uma civilização tecnologicamente
avançada.

Ao longo da campanha, Alex descobre que os ancestrais daquela
civilização deixaram a Terra milhões de anos atrás e que seu atual
planeta está entrando em colapso.

Uma parte do governo dos dinossauros pretende retornar à Terra,
mesmo que isso provoque uma guerra contra a humanidade.

Outra parte da população acredita que humanos e dinossauros podem
coexistir.

O jogador ficará no centro desse conflito.

\subsection{Gênero}

O jogo mistura os seguintes gêneros:

\begin{itemize}
    \item Aventura 3D;
    \item Ficção científica;
    \item Exploração;
    \item Sobrevivência;
    \item Furtividade;
    \item Puzzle;
    \item Ação leve;
    \item Narrativa.
\end{itemize}

A experiência principal será de aventura narrativa, com cada fase
introduzindo uma forma diferente de jogabilidade.

\subsection{Resumo do fluxo do jogo}

Na versão jogável atual:

\begin{center}
\textbf{Menu Principal}\\
$\downarrow$\\
\textbf{Nível 1 -- A Nave}\\
$\leftrightarrow$ \textbf{Pausa por ESC}\\
$\downarrow$\\
\textbf{Conclusão do Nível 1}\\
$\downarrow$\\
\textbf{Menu Principal}
\end{center}

Para a campanha completa, o planejamento mantém a sequência:

\begin{center}
Nave $\rightarrow$ Planeta $\rightarrow$ Laboratório\\
$\rightarrow$ Rebeldes $\rightarrow$ Terra $\rightarrow$ Final.
\end{center}

O ciclo implementado é explorar, identificar o objetivo, interagir,
resolver um obstáculo e liberar a próxima área. O HUD, as mensagens
contextuais, os indicadores e a sinalização orientam esse ciclo.

Diálogos, gravações e eventos narrativos encenados são recursos
previstos para a campanha, ainda não usados para conduzir a
progressão do Nível 1.



\subsection{Aparência e sensação}

O jogo será completamente 3D.

A direção visual combina ficção científica com elementos naturais e
pré-históricos.

O objetivo é fazer o jogador sentir que inicialmente está explorando
uma civilização alienígena, evitando revelar imediatamente que ela é
formada por dinossauros.

A estética da tecnologia utiliza:

\begin{itemize}
    \item metais escuros;
    \item luzes artificiais;
    \item painéis holográficos;
    \item estruturas circulares;
    \item símbolos desconhecidos;
    \item equipamentos de grande escala;
    \item arquitetura adaptada para criaturas maiores do que humanos.
\end{itemize}

As fases externas possuem vegetação densa, montanhas, rios, lagos,
pedras, cavernas e ruínas tecnológicas.

A câmera será em primeira pessoa para aumentar a sensação de
imersão e também manter o escopo do projeto controlado.

% =========================================================
% 3. HISTÓRIA, CENÁRIO E PERSONAGEM
% =========================================================

\section{História, Cenário e Personagem}

Esta seção apresenta o universo e o roteiro da campanha completa.
A história justifica o ponto de partida da nave, mas suas revelações,
personagens falantes e cenas posteriores ainda são planejamento.


\subsection{História e Narrativa}

\subsubsection*{Contexto}

Milhões de anos antes do surgimento da civilização humana, uma
pequena parcela de uma espécie inteligente originada dos dinossauros
desenvolveu tecnologia suficiente para detectar a aproximação de um
evento catastrófico na Terra.

Quando perceberam que a sobrevivência de sua civilização estava em
risco, construíram naves e abandonaram o planeta.

Eles encontraram um novo mundo habitável, posteriormente chamado de
\textit{Kryos}.

Durante milhões de anos, sua civilização se desenvolveu.

Kryos, porém, começou a apresentar um fenômeno desconhecido.

O núcleo do planeta tornou-se instável.

Terremotos, mudanças climáticas e falhas geológicas passaram a
destruir regiões inteiras.

Os cientistas determinaram que Kryos eventualmente se tornaria
inabitável.

A descoberta causou uma divisão política e social.

O governo criou o \textbf{Projeto Retorno}.

Seu objetivo era retornar ao planeta de origem: a Terra.

Quando as sondas descobriram que uma nova espécie inteligente havia
dominado o planeta, o problema tornou-se mais complexo.

Essa espécie era a humanidade.

O governo iniciou secretamente o \textbf{Programa Observador}.

Humanos passaram a ser sequestrados para estudos biológicos,
comportamentais e ambientais.

O objetivo era descobrir:

\begin{itemize}
    \item se os dinossauros poderiam sobreviver na Terra moderna;
    \item quais doenças humanas poderiam afetá-los;
    \item qual era a capacidade militar humana;
    \item como humanos reagiriam ao contato com outra civilização;
    \item se uma possível coexistência seria viável.
\end{itemize}

Alex Moraes é um dos humanos selecionados.

O jogo começa quando uma falha no sistema da nave interrompe seu
estado de contenção.

\subsubsection*{Progressão narrativa}

A história é revelada de maneira gradual.

\begin{enumerate}
    \item Alex acredita ter sido sequestrado por alienígenas.
    \item Alex consegue escapar de sua câmara.
    \item Ele percebe que está dentro de uma nave de pesquisa desconhecida.
    \item Ao sair da nave, descobre que está em outro planeta.
    \item Ele encontra pegadas gigantes.
    \item Escuta rugidos e observa grandes sombras.
    \item Finalmente encontra um dinossauro inteligente.
    \item Descobre que os dinossauros conhecem a humanidade.
    \item Descobre que humanos estão sendo estudados.
    \item Descobre que Kryos está morrendo.
    \item Descobre a existência do Projeto Retorno.
    \item Conhece um grupo rebelde.
    \item Descobre que os ancestrais dos dinossauros vieram da Terra.
    \item Descobre que parte do governo pretende tomar a Terra pela força.
    \item Retorna à Terra.
    \item Tenta impedir uma guerra entre as duas espécies.
\end{enumerate}

\subsubsection*{Cenas de diálogo principais -- roteiro planejado}

As cenas e falas desta subseção pertencem ao planejamento narrativo.
A versão atual não possui diálogos ou cutscenes implementados.


\paragraph{Cena 1 -- Despertar}

Na implementação atual, o controle começa diretamente na sala de
contenção; a encenação abaixo permanece prevista para outra etapa.


Alex acorda dentro da câmara.

Uma voz automática reproduz uma mensagem incompreensível.

Após alguns segundos, o sistema apresenta uma tradução incompleta:

\begin{quote}
``Organismo 31... consciência detectada... contenção comprometida.''
\end{quote}

Alex responde:

\begin{quote}
``Onde eu estou?''
\end{quote}

Não há resposta.

\paragraph{Cena 2 -- Primeiro encontro}

Na segunda fase, Alex observa uma criatura caminhando entre árvores.

Inicialmente apenas sua sombra é visível.

A criatura se aproxima.

Alex percebe que se trata de um dinossauro utilizando equipamentos
tecnológicos.

Alex:

\begin{quote}
``Isso não é possível.''
\end{quote}

\paragraph{Cena 3 -- Laboratório}

Alex encontra registros sobre humanos.

Uma transmissão traduzida informa:

\begin{quote}
``Espécie dominante do planeta de origem confirmada.
Compatibilidade ambiental: elevada.''
\end{quote}

Alex percebe que a Terra é o planeta citado.

\paragraph{Cena 4 -- Rebeldes}

A líder rebelde Tala explica:

\begin{quote}
``Vocês chamam nosso planeta de Terra. Nós o chamávamos de origem.''
\end{quote}

Nesse momento, Alex compreende que aquela civilização não encontrou a
Terra recentemente.

Eles nasceram nela.

\paragraph{Cena 5 -- Confronto final}

O comandante Rhaz afirma:

\begin{quote}
``Kryos está morrendo. Não condenarei nossa espécie por uma civilização
que surgiu depois de nossa partida.''
\end{quote}

Alex responde:

\begin{quote}
``Sobreviver não precisa significar destruir quem já está lá.''
\end{quote}

% =========================================================
% LENTE 1
% =========================================================

\subsubsection*{Lente \#1: Experiência Essencial}

\paragraph{Quais experiências queremos que o jogador tenha?}

O jogador deve experimentar:

\begin{itemize}
    \item mistério;
    \item curiosidade;
    \item sensação de isolamento;
    \item descoberta;
    \item medo do desconhecido;
    \item surpresa;
    \item perseguição;
    \item exploração;
    \item tensão;
    \item empatia por personagens inicialmente considerados inimigos;
    \item conflito moral.
\end{itemize}

A experiência principal é fazer o jogador mudar constantemente sua
interpretação sobre a situação.

No início, o pensamento esperado é:

\begin{quote}
``Fui sequestrado por alienígenas e preciso fugir.''
\end{quote}

Mais tarde:

\begin{quote}
``Os alienígenas são dinossauros.''
\end{quote}

Depois:

\begin{quote}
``Eles conhecem a Terra.''
\end{quote}

Em seguida:

\begin{quote}
``Eles vieram da Terra.''
\end{quote}

E finalmente:

\begin{quote}
``Nem todos são meus inimigos.''
\end{quote}

\paragraph{O que é essencial para essa experiência?}

São considerados eventos essenciais:

\begin{enumerate}
    \item Despertar sem explicações.
    \item Explorar a nave.
    \item Observar criaturas sem enxergá-las completamente.
    \item Escapar da nave.
    \item Descobrir o planeta.
    \item Encontrar as primeiras pegadas.
    \item Encontrar o primeiro dinossauro.
    \item Ser levado ao laboratório.
    \item Encontrar registros de humanos.
    \item Descobrir o Projeto Retorno.
    \item Sofrer uma queda durante a tentativa de fuga.
    \item Conhecer os rebeldes.
    \item Descobrir a origem terrestre dos dinossauros.
    \item Retornar para a Terra.
    \item Impedir o conflito total.
\end{enumerate}

\paragraph{Como o jogo captura essa essência?}

O mistério será representado pela ausência proposital de informações
no início.

O jogo evita grandes explicações nas primeiras fases.

O jogador encontra pequenos fragmentos de informação através de
cenários, gravações, computadores e diálogos.

As revelações são apresentadas principalmente como consequência da
exploração, e não apenas através de textos.

Pegadas, ruídos, sombras, equipamentos e arquitetura funcionam como
pistas visuais.

% =========================================================
% LENTE 2
% =========================================================

\subsubsection*{Lente \#2: Surpresa}

Os principais momentos de surpresa serão:

\begin{enumerate}
    \item A descoberta de que Alex não está na Terra.
    \item As primeiras pegadas gigantes.
    \item O primeiro dinossauro inteligente.
    \item A descoberta de humanos presos no laboratório.
    \item A descoberta de que os dinossauros conhecem a Terra.
    \item A revelação de que Kryos está morrendo.
    \item O encontro com dinossauros amigáveis.
    \item A revelação de que os ancestrais dos dinossauros vieram da Terra.
    \item A descoberta de que a invasão ainda pode ser evitada.
\end{enumerate}

As regras também permitem pequenas surpresas dinâmicas.

Por exemplo:

\begin{itemize}
    \item portas podem liberar novos caminhos;
    \item inimigos podem alterar suas rotas;
    \item alarmes podem modificar o comportamento dos NPCs;
    \item algumas salas são opcionais;
    \item documentos opcionais revelam informações extras;
    \item determinados acontecimentos podem ser vistos apenas por jogadores
    que explorarem determinadas áreas.
\end{itemize}

% =========================================================
% 3.2 MUNDO
% =========================================================

\subsection{Mundo do Jogo}

\subsubsection{Visual geral e sensação do mundo}

O mundo mistura elementos pré-históricos e futuristas.

A direção narrativa da campanha prevê que a tecnologia dos dinossauros
pareça desconhecida ao humano. O Nível 1 atual utiliza kits sci-fi
modulares e escala transitável por um personagem humano.

As construções possuem:

\begin{itemize}
    \item portas grandes;
    \item corredores amplos;
    \item terminais elevados;
    \item plataformas;
    \item grandes estruturas metálicas;
    \item iluminação artificial;
    \item símbolos próprios;
    \item hologramas;
    \item estruturas adaptadas para diferentes espécies.
\end{itemize}

O planeta Kryos possui grandes regiões naturais preservadas.

Isso cria contraste entre tecnologia extremamente avançada e natureza
selvagem.

\subsubsection{Áreas do mundo}

\paragraph{Nave de Pesquisa}

Utilizada na Fase 1.

É uma nave de pesquisa responsável pelo transporte de indivíduos
capturados e equipamentos científicos. O cenário jogável atual tem
porte médio e nove setores, com circulação em primeira pessoa.

Possui:

\begin{itemize}
    \item câmaras de contenção;
    \item corredores;
    \item laboratório;
    \item depósitos;
    \item sala de energia;
    \item centro de segurança;
    \item hangares.
\end{itemize}

\paragraph{Vale de Kryos}

Utilizado na Fase 2.

É uma região natural próxima da nave.

Possui:

\begin{itemize}
    \item florestas;
    \item rios;
    \item morros;
    \item vales;
    \item pedras;
    \item areia;
    \item cascalho;
    \item cavernas;
    \item ruínas;
    \item estruturas tecnológicas.
\end{itemize}

\paragraph{Complexo de Pesquisa}

Utilizado na Fase 3.

É o local onde espécies terrestres são estudadas.

Possui:

\begin{itemize}
    \item laboratórios;
    \item celas;
    \item corredores;
    \item câmeras;
    \item salas de controle;
    \item áreas médicas;
    \item centros de contenção.
\end{itemize}

\paragraph{Território Rebelde}

Utilizado na Fase 4.

É uma região distante das grandes cidades.

Possui:

\begin{itemize}
    \item florestas;
    \item cânions;
    \item destroços;
    \item acampamentos;
    \item cavernas;
    \item equipamentos improvisados;
    \item antiga instalação de comunicação.
\end{itemize}

\paragraph{Terra}

Utilizada na Fase 5.

O jogador retorna a uma região urbana da Terra.

A área foi parcialmente evacuada devido à chegada das primeiras naves.

Possui:

\begin{itemize}
    \item ruas;
    \item casas;
    \item prédios;
    \item veículos;
    \item barricadas;
    \item destroços;
    \item postos militares;
    \item estruturas de comunicação.
\end{itemize}

% =========================================================
% 3.3 PERSONAGENS
% =========================================================

\subsection{Personagens}

\subsubsection*{Alex Moraes}

\textbf{Espécie:} Humano.

\textbf{Idade:} 28 anos.

\textbf{História:}

Alex trabalhava como técnico de manutenção eletrônica antes de ser
capturado.

Essa profissão justifica sua capacidade de interagir com máquinas,
circuitos e terminais durante o jogo.

\textbf{Personalidade:}

\begin{itemize}
    \item curioso;
    \item observador;
    \item desconfiado;
    \item persistente;
    \item inicialmente impulsivo;
    \item progressivamente mais empático.
\end{itemize}

\textbf{Aparência:}

Roupa simples modificada durante o período de contenção.
Posteriormente utiliza equipamentos encontrados em Kryos.

\textbf{Situação atual:} o personagem é controlado em primeira pessoa.

Caminhar, correr, saltar e agachar são ações funcionais, mas não há
modelo humano com todos os estados animados comprovados.

\textbf{Animações planejadas para o protagonista:}

\begin{itemize}
    \item parado;
    \item queda;
    \item caminhada;
    \item corrida;
    \item salto;
    \item interação;
    \item dano;
    \item morte.
\end{itemize}

A representação futura poderá utilizar braços ou um modelo humano.
Movimentar a câmera ou animar o pet não é tratado como comprovação
automática dos estados de animação exigidos para o personagem.

\textbf{Habilidades:}

\begin{itemize}
    \item correr;
    \item saltar;
    \item agachar;
    \item interagir;
    \item utilizar terminais;
    \item carregar objetos pequenos;
    \item utilizar equipamentos;
    \item usar arma simples nas fases finais.
\end{itemize}

\textbf{Relevância:}

Protagonista e ponto de vista principal da narrativa.

% ---------------------------------------------------------

\subsubsection*{Tala}

\textbf{Espécie:} Descendente inteligente de ceratopsídeos.

\textbf{Função:} Líder de uma célula rebelde.

\textbf{Personalidade:}

\begin{itemize}
    \item racional;
    \item paciente;
    \item protetora;
    \item desconfiada do governo;
    \item favorável à coexistência.
\end{itemize}

\textbf{Aparência:}

Possui equipamentos improvisados, proteção corporal e acessórios
tecnológicos reutilizados.

\textbf{Animações previstas:}

\begin{itemize}
    \item idle;
    \item andar;
    \item correr;
    \item conversar;
    \item apontar;
    \item atacar;
\end{itemize}

\textbf{Habilidades:}

Conhecimento do planeta, liderança e combate.

\textbf{Relevância:}

É responsável por explicar a verdadeira história da civilização de
Kryos.

\textbf{Relacionamento com Alex:}

Inicialmente desconfia do protagonista, mas passa a considerá-lo uma
possível ligação entre as duas espécies.

% ---------------------------------------------------------

\subsubsection*{General Rhaz}

\textbf{Espécie:} Descendente de grande terópode.

\textbf{Função:} Comandante responsável pelo Projeto Retorno.

\textbf{Personalidade:}

\begin{itemize}
    \item autoritário;
    \item disciplinado;
    \item pragmático;
    \item extremamente focado na sobrevivência de sua espécie.
\end{itemize}

Rhaz não considera a destruição dos humanos seu objetivo principal.

Para ele, uma possível guerra é uma consequência aceitável caso seja
necessária para salvar seu povo.

\textbf{Aparência:}

Grande porte, armadura tecnológica e dispositivos de comunicação.

\textbf{Relevância:}

Principal antagonista da campanha.

% ---------------------------------------------------------

\subsubsection*{Dra. Veyra}

\textbf{Espécie:} Descendente de pequeno terópode.

\textbf{Função:} Cientista responsável pelo estudo de humanos.

\textbf{Personalidade:}

\begin{itemize}
    \item lógica;
    \item curiosa;
    \item fria;
    \item cientificamente interessada nos humanos;
    \item gradualmente questiona as ordens recebidas.
\end{itemize}

\textbf{Relevância:}

É responsável por parte dos registros que permitem ao jogador
compreender o Programa Observador.

% =========================================================
% 4. JOGABILIDADE
% =========================================================

\subsubsection*{Sparrow -- Companheiro opcional}

\textbf{Representação:} pequeno pássaro do pacote Sparrow --
Quirky Series, da Omabuarts Studio.

\textbf{Função atual:} encontro opcional depois da primeira porta.
Permanece no chão até o jogador fazer amizade com F e, então,
acompanha o protagonista em voo.

\textbf{Personalidade visual:} repouso variado, pequenos movimentos
corporais do próprio pacote e voo suave. Não possui falas,
combate, coleta de itens ou participação obrigatória nos puzzles.

\textbf{Animações utilizadas:} Idle\_A, Idle\_B, Idle\_C e Fly.

\textbf{Áudio:} componente de cantos ocasionais preparado, sem clipes
adequados atribuídos na versão atual.

\textbf{Persistência:} adoção somente durante a partida; reiniciar
faz o pássaro voltar ao ponto de espera.


\section{Jogabilidade e Mecânicas}

\subsection{Jogabilidade / Gameplay}

\subsubsection{Objetivos}

O objetivo geral é sobreviver, compreender o motivo do sequestro e
retornar à Terra.

Os objetivos mudam conforme a fase.

\begin{enumerate}
    \item Escapar da nave.
    \item Explorar o planeta e encontrar abrigo.
    \item Escapar do laboratório.
    \item Ajudar os rebeldes e conseguir uma forma de retornar.
    \item Impedir o início de uma guerra na Terra.
\end{enumerate}

\subsubsection{Progressão do Jogo}

A campanha planejada possui cinco fases principais; somente a primeira
está implementada na versão atual.

Cada fase introduz ou destaca uma mecânica.

\begin{center}
\begin{tabularx}{\textwidth}{|c|X|X|}
\hline
\textbf{Fase} & \textbf{Ambiente} & \textbf{Gameplay principal}\\
\hline
1 & Nave & Exploração, arremesso e puzzles\\
\hline
2 & Planeta & Exploração e sobrevivência\\
\hline
3 & Laboratório & Furtividade\\
\hline
4 & Território Rebelde & Exploração e fuga\\
\hline
5 & Terra & Ação e objetivo narrativo\\
\hline
\end{tabularx}
\end{center}

\subsubsection{Estrutura da Missão/Desafio}

O caminho obrigatório implementado no Nível 1 é:

\begin{enumerate}
\item Pegar uma caixa na sala de contenção.
\item Arremessar a caixa contra o alvo acima da porta inicial.
\item Encontrar a bateria em Máquinas/Carga, acessível pelo lado
esquerdo do depósito.
\item Instalar a bateria no painel da sala de energia.
\item Entrar no laboratório desbloqueado e encontrar o cartão.
\item Usar o cartão no leitor da segurança para liberar o hangar.
\item Ativar os três terminais do hangar, em qualquer ordem.
\item Atravessar a saída iluminada para concluir o nível.
\end{enumerate}

Adotar o Sparrow é uma atividade opcional depois da primeira porta
e não bloqueia a progressão.

\subsubsection{Estrutura dos Enigmas}

Os desafios utilizam exploração, observação e interação simples.
Não há senhas, ordem secreta de terminais ou minigames na versão atual.

\paragraph{Impacto} Pegar a caixa com F e lançá-la com o botão esquerdo
do mouse contra o alvo físico acima da porta. Erros permitem novas tentativas.

\paragraph{Energia} Coletar a bateria e instalá-la no painel. A energia
restaurada libera o laboratório.

\paragraph{Acesso} Encontrar o único cartão, sorteado entre três pontos
válidos do laboratório, e utilizá-lo no leitor da segurança.

\paragraph{Três sistemas} Ativar uma vez cada terminal do hangar.
A ordem é livre. Somente o conjunto completo libera a saída externa.



\subsubsection{Fluxo do jogo}

O fluxo geral é:

\begin{center}
Exploração
$\rightarrow$
Objetivo
$\rightarrow$
Obstáculo
$\rightarrow$
Interação/Puzzle
$\rightarrow$
Recompensa
$\rightarrow$
Nova área.
\end{center}

Na campanha planejada, as recompensas poderão incluir os itens abaixo.
No Nível 1 atual, a recompensa funcional é principalmente o acesso
à próxima área e a atualização do objetivo; checkpoints e documentos
narrativos ainda não estão implementados.

Possibilidades da campanha:

\begin{itemize}
    \item acesso a uma nova região;
    \item informação narrativa;
    \item equipamento;
    \item checkpoint;
    \item novo objetivo.
\end{itemize}

% =========================================================
% MECÂNICAS
% =========================================================

\subsection{Mecânicas}

\subsubsection{Física e colisões}

O jogador usa o \texttt{CharacterController} existente, gravidade e
colisão com o cenário. Não utiliza Rigidbody para a movimentação.

As caixas arremessadas usam Rigidbody, impulso com \texttt{AddForce},
torque com \texttt{AddTorque} e detecção contínua de colisão. A direção
do lançamento acompanha a câmera.

Pisos, paredes, tetos, portas fechadas e equipamentos sólidos possuem
colisores adequados à geometria. São priorizados BoxColliders e
composições simples. Mesas usam tampo e pés; marcos de portas usam
laterais e verga. Cabos baixos não se tornam barreiras desnecessárias.

Cada porta de progressão desloca sua folha por \texttt{Translate}.
O collider sólido acompanha a folha e libera o vão ao abrir.
Triggers de interação e conclusão são separados da colisão física.

\subsubsection{Movimento no jogo}

\begin{itemize}
\item WASD: andar.
\item Mouse: olhar ao redor, com limite vertical.
\item Shift esquerdo: correr.
\item Espaço: saltar.
\item Ctrl esquerdo ou direito, mantido pressionado: agachar.
\end{itemize}

O agachamento reduz a cápsula e a altura da câmera sem deslocar os
pés. O jogador só fica em pé quando há espaço acima. Não há stamina.

\subsubsection{Interação em primeira pessoa}

\texttt{PlayerInteraction} utiliza Raycast a partir da câmera, com
alcance máximo de aproximadamente 3 metros. A tecla principal é
\textbf{F}. A mensagem contextual aparece na parte inferior da tela.

Obstáculos sólidos interrompem a busca. Não se coletam itens ou se
acionam painéis através de paredes ou portas fechadas.
\texttt{IInteractable} é o contrato compartilhado pelos interativos
e pela adoção do Sparrow. A tecla E não é utilizada nesta versão.

Interações e lançamentos são bloqueados na pausa e após a conclusão.

\subsubsection{Arremesso e spawn de caixas}

A caixa inicial fica no chão da contenção. Após F, uma instância
aparece no \texttt{HoldPoint}, filho da câmera, abaixo e ao lado da
mira. Enquanto segura, não colide com o jogador nem intercepta seus
Raycasts.

O botão esquerdo lança uma caixa física. A função de spawn usa
\texttt{Instantiate} para criar automaticamente outra na mão.
Um array contém dois prefabs locais de caixas baseados em modelos
reais, selecionados com \texttt{Random.Range}.

Os parâmetros atuais são impulso 11, torque 0,08, intervalo mínimo
de 0,45 segundo, duração de 12 segundos e limite de oito projéteis
ativos. São configuráveis no Inspector.

Somente uma caixa efetivamente lançada pelo sistema ativa o alvo.
O acerto abre a primeira porta, confirma visualmente o sucesso,
remove as caixas e interrompe a geração de reposições.

\subsubsection{Itens e estado dos puzzles}

Bateria e cartão são itens-chave registrados por estado, sem
inventário visual complexo. A coleta remove o item do cenário.
A bateria instalada aparece no encaixe do painel.

O cartão reutiliza o modelo ARCADIA, com \texttt{Arcadia\_Verso}
no verso. Há somente um cartão ativo. No início de cada partida,
é sorteado em execução entre três pontos apoiados em mesas reais
do laboratório. São verificados limites, apoio e espaço livre.

\texttt{Level01GameManager} é a única fonte de verdade para energia,
cartão, terminais e conclusão, evitando estados conflitantes.

\subsubsection{Passarinho companheiro}

O Sparrow é opcional. Começa no chão de um ponto visível do corredor
logo depois da primeira porta, sem acompanhar o jogador.

Ao mirar nele e pressionar F, aparece a mensagem
\emph{Novo companheiro encontrado!}. A adoção é registrada apenas na
partida. O alvo de interação é desativado e o pássaro decola para
seguir suavemente ao lado do protagonista.

O movimento usa interpolação, varreduras contra obstáculos e uma
trilha curta do jogador. Se necessário, a recuperação ocorre de
forma oculta em posição segura fora da vista frontal. Não usa NavMesh.

Após a adoção, o pet não tem colisão física ativa, não bloqueia
Raycasts, não aciona puzzles ou conclusão e não interfere nas caixas.
Pausa com ESC e volta a esperar ao reiniciar.

\subsubsection{Combate e furtividade}

Não há combate, inimigos, dano, vida, drone ou furtividade funcional
no Nível 1 atual. Câmeras de vigilância são decorativas.

A campanha prevê furtividade no Nível 3 e ação simples nas etapas
finais. São sistemas futuros, ainda não implementados.

\subsubsection{Economia}

Não há dinheiro, comércio ou loja. A progressão depende de itens-chave,
ações concluídas e acesso a novas áreas.

\subsubsection{Fluxo de telas atual}

\begin{center}
\textbf{MainMenu: Novo Jogo / Controles / Sair}\\[0.2cm]
$\downarrow$\\[0.2cm]
\textbf{Level01\_Nave\_Expanded}\\[0.2cm]
$\leftrightarrow$ \textbf{Pausa por ESC}\\[0.2cm]
$\downarrow$\\[0.2cm]
\textbf{NÍVEL 1 CONCLUÍDO}\\[0.2cm]
$\downarrow$\\[0.2cm]
\textbf{Voltar ao Menu Principal}
\end{center}

A campanha completa mantém cinco níveis planejados e um final
narrativo. Essa sequência ainda não é jogável na versão atual.

\subsection{Opções do jogo}

O menu oferece a consulta dos controles. Não há tela de ajustes de
volume, sensibilidade, resolução ou qualidade gráfica. A equipe
pode editar parâmetros no Inspector.

\subsection{Rejogabilidade e salvamento}

Não há save, carregamento de progresso ou checkpoints. Novo Jogo e
Reiniciar Nível recarregam a cena e zeram os estados.

A porta inicial fecha, caixas antigas desaparecem, bateria retorna,
cartão é sorteado novamente, terminais são resetados, saída e luz
voltam ao início, e o Sparrow espera sem adoção.

O sorteio do cartão e as variações de idle acrescentam pequena
variedade sem alterar a solução principal.

\subsection{Exploração opcional e Easter Eggs}

Explorar salas auxiliares e adotar o Sparrow são atividades opcionais.
Não há sala secreta, senha especial, documento coletável ou Easter Egg
funcional declarado como implementado. Referências adicionais poderão
ser incluídas depois, sem complicar o percurso obrigatório.



\section{Níveis}

\subsection{Nível 1 -- A Nave}

\subsubsection*{Sinopse e estado atual}

Alex inicia na sala de contenção de uma nave de pesquisa desconhecida.
A arquitetura e os equipamentos sugerem uma origem alienígena.
O jogador precisa liberar sua primeira passagem, restaurar energia
e abrir a saída externa do hangar.

A fase está jogável do início ao fim. O controle é entregue
diretamente na contenção, com o objetivo de pegar uma caixa.
Não há cutscene de despertar, abertura animada de cápsula,
voz de apresentação ou inimigo nesta versão.

A nave é de tamanho médio, modular e explorável. O objetivo é
oferecer ambientes distintos e caminhos legíveis, sem alongar
artificialmente o percurso. Uma duração futura de 5 a 10 minutos
é uma intenção de design, não uma medição validada da partida atual.

\subsubsection*{Setores e função}

\begin{center}
\begin{tabularx}{\textwidth}{|p{3.5cm}|X|}
\hline
\textbf{Setor} & \textbf{Função na versão atual}\\
\hline
01 -- Contenção & Spawn, caixa inicial, alvo de impacto e primeira porta bloqueada.\\
\hline
02 -- Corredores & Conexão entre setores, sinalização e Sparrow esperando depois da primeira porta.\\
\hline
03 -- Energia & Painel acessível desde o início; instalação da bateria libera o laboratório.\\
\hline
04 -- Depósito & Armazenamento e acesso ao desvio de Máquinas/Carga.\\
\hline
05 -- Laboratório & Porta dependente da energia; três locais possíveis para o único cartão.\\
\hline
06 -- Segurança & Leitor do cartão que autoriza o acesso ao hangar.\\
\hline
07 -- Hangar & Três terminais, portão externo, luz branca e região de conclusão.\\
\hline
08 -- Máquinas/Carga & Ambiente industrial à esquerda do depósito, onde fica a bateria.\\
\hline
09 -- Manutenção & Sala auxiliar para exploração opcional.\\
\hline
\end{tabularx}
\end{center}

Os corredores conectam os setores. Visitar energia antes de possuir
a bateria é permitido, mas não resolve o puzzle. A bateria é acessível
antes da restauração; o cartão está no laboratório, cuja abertura
depende somente da bateria. Não há dependência circular.

\subsubsection*{Objetivos apresentados pelo HUD}

\begin{enumerate}
\item Pegue uma caixa na sala inicial.
\item Acerte o botão acima da porta com a caixa.
\item Encontre uma bateria na nave.
\item Instale a bateria na sala de energia.
\item Entre no laboratório.
\item Encontre o cartão de segurança.
\item Utilize o cartão no leitor da segurança.
\item Ative os três terminais do hangar, com contador de 0/3 a 3/3.
\item Atravesse a saída para escapar da nave.
\end{enumerate}

A adoção não substitui o objetivo principal. A confirmação do
companheiro é uma mensagem temporária.

\subsubsection*{Desafio inicial -- Impacto acima da porta}

A porta da contenção começa fechada e sólida. O alvo fica acima dela,
visível da sala. Uma caixa pequena permite iniciar o desafio com F.

O jogador arremessa com o mouse esquerdo, observando a direção e a
gravidade. Após cada lançamento, uma nova caixa aparece na mão.
Um intervalo curto e o limite de projéteis evitam acúmulo excessivo.

Um impacto válido da caixa lançada muda o indicador do alvo e abre
a porta. Contato do personagem, do pet ou de decoração não é aceito.
Após o sucesso, o sistema interrompe o spawn e remove as caixas.

\subsubsection*{Exploração opcional -- Adoção do Sparrow}

Depois da porta inicial, o jogador encontra o Sparrow no chão do
corredor, próximo à lateral e fora da passagem principal. Ele utiliza
os idles reais do pacote enquanto espera.

A mensagem é \textbf{[F] Fazer amizade com o passarinho}. Ao adotar,
o pet passa a acompanhar o jogador em voo. A adoção pode ser ignorada
e não fornece bateria, cartão ou vantagem obrigatória nos puzzles.

\subsubsection*{Puzzle 1 -- Bateria e energia}

A bateria utiliza uma célula real do pacote \textit{Energy Cells Free}.
Fica em Máquinas/Carga, no desvio à esquerda do depósito.

Ao coletar com F, desaparece do cenário, o estado de posse é
registrado e aparece \emph{Bateria coletada!}. O objetivo passa a
indicar a instalação na sala de energia.

O painel sem bateria informa \emph{Bateria necessária}. Com o item,
F o instala, mostra a bateria no encaixe, ativa o indicador, reproduz
a confirmação sonora e o VFX ciano. A porta do laboratório abre.

\subsubsection*{Puzzle 2 -- Cartão de segurança}

Dentro do laboratório há um cartão ARCADIA. Sua posição é sorteada
entre três marcadores válidos sobre mesas, em cada início ou reinício.
O modelo não é duplicado.

O jogador mira e pressiona F. O cartão desaparece, a posse é
registrada e o HUD orienta a ida ao leitor da segurança.

Sem cartão, o leitor informa \emph{Cartão de acesso necessário}.
Com o cartão, a interação autoriza o acesso e abre a porta do hangar.
Não existe minigame de autenticação.

\subsubsection*{Puzzle 3 -- Sistemas do hangar}

Os três terminais representam:

\begin{enumerate}
\item Sistema de energia auxiliar.
\item Sistema de segurança.
\item Controle da saída externa.
\end{enumerate}

Cada terminal é ativado uma vez com F e muda seu indicador.
A ordem é livre; não existe senha ou sequência de botões.
O HUD apresenta o progresso de 0/3 até 3/3.

Com os três ativos, o portão externo abre, o vapor é emitido e um
alerta eletrônico toca por quatro segundos. A luz branca fica
visível atrás do portão. Repetir uma solicitação não duplica o alerta.

\subsubsection*{Estados das portas}

\begin{center}
\begin{tabularx}{\textwidth}{|p{3.3cm}|X|}
\hline
\textbf{Porta inicialmente fechada} & \textbf{Condição de abertura}\\
\hline
Contenção & Impacto válido de caixa lançada no alvo.\\
\hline
Laboratório & Bateria instalada e energia restaurada.\\
\hline
Entrada do hangar & Cartão utilizado no leitor da segurança.\\
\hline
Saída externa & Três terminais do hangar ativados.\\
\hline
\end{tabularx}
\end{center}

Todas bloqueiam fisicamente quando fechadas. A colisão acompanha a
folha móvel ao abrir.

\subsubsection*{Conclusão}

Abrir o portão não conclui sozinho a fase. O jogador deve caminhar
até a luz branca e entrar no volume de saída.

O trigger identifica somente a cápsula do jogador, verifica sua
presença no volume e solicita a conclusão ao gerente. São exigidos
os três terminais e a liberação do portão. Caixas e pássaro não contam.

A tela apresenta:

\begin{center}
\textbf{NÍVEL 1 CONCLUÍDO}\\[0.15cm]
Você conseguiu escapar da nave.\\[0.15cm]
\textbf{Voltar ao Menu Principal}
\end{center}

O jogo congela, libera o cursor e permite voltar ao menu.
Não carrega automaticamente o Nível 2 nem mostra um planeta
explorável ou a revelação dos dinossauros.

\subsubsection*{Apresentação audiovisual}

A nave combina iluminação fria, painéis, computadores, caixas,
tubulações e equipamentos industriais. Os setores mantêm identidade
própria sem prejudicar a circulação.

As quatro portas importantes usam som mecânico 3D e pequenas
descargas de vapor. A bateria instalada ativa faíscas ciano.
Os efeitos são curtos, sem loop e sem colisores.

O áudio de ambiente é discreto na energia. Passos e cantos do
Sparrow têm componentes preparados, mas ainda não possuem clipes
adequados atribuídos. Não são declarados como sons audíveis atuais.

\subsubsection*{Assets e comprovação da variedade}

O inventário formal da expansão registra \textbf{82 modelos visuais
distintos em 1005 instâncias de cenário}. Esse é o catálogo-base
auditado, anterior aos últimos acréscimos de gameplay e polimento,
e não uma nova contagem total da cena atual.

O critério agrupa geometrias equivalentes: repetir prefab, trocar
material, mudar escala ou rotação não cria um novo modelo.
UI, materiais, áudios, objetos vazios e repetições não aumentam
o total de modelos visuais.

A distribuição do catálogo-base é:

\begin{center}
\begin{tabularx}{\textwidth}{|X|r|}
\hline
\textbf{Categoria} & \textbf{Modelos distintos}\\
\hline
Armazenamento & 20\\
Arquitetura & 21\\
Científico & 2\\
Computadores & 5\\
Estruturas industriais & 6\\
Iluminação & 5\\
Industrial & 14\\
Mobiliário & 6\\
Portas & 2\\
Segurança & 1\\
\hline
\textbf{Total da base auditada} & \textbf{82}\\
\hline
\end{tabularx}
\end{center}

As origens, setores e quantidades estão no relatório
\texttt{Docs/Relatorios/Nivel01\_Expandido/Level01\_Expanded\_Assets.md}
e nos CSV/JSON correspondentes. A contagem final deve ser
regenerada a partir da cena entregue, sem somar estimativas
aos modelos comprovados.

Pacotes utilizados na nave incluem Barking Dog -- 3D Free Modular
Kit, Modular SciFi MegaKit, Free Sci-Fi Office Pack, Low Poly
Storage Pack e Surveillance Camera. Os itens usam Energy Cells
Free e o cartão ARCADIA já existente. O companheiro utiliza
Sparrow -- Quirky Series, da Omabuarts Studio.

\subsubsection*{Limites desta entrega}

Não estão implementados drone, inimigos, perseguição, combate,
vida, inventário complexo, save, checkpoints, diálogos,
cutscenes, abertura encenada da cápsula ou transição para outro nível.
As câmeras são decorativas. Esses elementos não são necessários
para resolver os puzzles atuais.



\subsection{Nível 2 -- O Planeta}

\textbf{Status: planejado, não implementado na versão atual.}


\subsubsection*{Sinopse}

Alex deixa a nave e começa a explorar Kryos.

Não existe caminho evidente para casa.

O jogador precisa atravessar uma região natural em busca de uma
estrutura observada à distância.

Durante a exploração, começa a encontrar evidências de criaturas
gigantes.

\subsubsection*{Introdução}

A fase começa imediatamente após a saída da nave.

O jogador vê:

\begin{itemize}
    \item montanhas;
    \item árvores;
    \item grandes luas;
    \item rios;
    \item estruturas artificiais ao longe.
\end{itemize}

Objetivo:

\begin{center}
\textbf{Encontre uma forma de pedir ajuda.}
\end{center}

\subsubsection*{Objetivos}

\begin{enumerate}
    \item Descer da região da nave.
    \item Seguir sinais de uma instalação.
    \item Atravessar a floresta.
    \item Investigar pegadas.
    \item Encontrar abrigo.
    \item Atravessar um rio ou vale.
    \item Alcançar uma estrutura.
    \item Encontrar o primeiro dinossauro.
\end{enumerate}

\subsubsection*{Eventos}

Primeiro, o jogador encontra uma pegada.

Depois escuta um rugido.

Mais tarde vê árvores se movimentando.

Uma grande sombra passa ao longe.

Somente depois disso ocorre a revelação.

Uma criatura aparece.

Ela possui roupas, equipamentos e tecnologia.

Alex percebe que é um dinossauro inteligente.

\subsubsection*{Mecânicas}

\begin{itemize}
    \item exploração de terreno;
    \item água;
    \item terreno irregular;
    \item checkpoints;
    \item coleta;
    \item áreas de perigo;
    \item partículas ambientais;
    \item eventos ativados por Trigger.
\end{itemize}

\subsubsection*{Ambiente}

O terreno deve utilizar:

\begin{itemize}
    \item grama;
    \item terra;
    \item areia;
    \item cascalho;
    \item pedra;
    \item água;
    \item árvores;
    \item arbustos;
    \item morros;
    \item vales.
\end{itemize}

\subsubsection*{Meta de distribuição de assets -- planejamento}

Os números abaixo são metas de modelos visuais distintos para o
nível futuro, não um inventário de assets importados ou posicionados.
Repetições, mudanças de material e texturas isoladas não aumentam
essa contagem. A distribuição precisa ser comprovada na implementação.


\begin{center}
\begin{tabularx}{\textwidth}{|X|c|}
\hline
Vegetação & 24\\
\hline
Pedras e formações naturais & 15\\
\hline
Estruturas de trilha e passagem & 8\\
\hline
Ruínas e estruturas & 10\\
\hline
Objetos tecnológicos & 7\\
\hline
Elementos de fauna & 5\\
\hline
Elementos especiais & 5\\
\hline
\textbf{Total previsto} & \textbf{74}\\
\hline
\end{tabularx}
\end{center}

% =========================================================
% FASE 3
% =========================================================

\subsection{Nível 3 -- O Laboratório}

\textbf{Status: planejado, não implementado na versão atual.}


\subsubsection*{Sinopse}

Alex é capturado novamente.

Quando acorda, está dentro de um centro de pesquisa.

Ali descobre que não é o primeiro humano levado para Kryos.

A fase é focada em furtividade.

\subsubsection*{Objetivos}

\begin{enumerate}
    \item Escapar da cela.
    \item Evitar os guardas.
    \item Encontrar o centro de segurança.
    \item Desativar câmeras.
    \item Encontrar registros.
    \item Descobrir o propósito dos sequestros.
    \item Libertar outros prisioneiros, opcionalmente.
    \item Encontrar uma rota de fuga.
\end{enumerate}

\subsubsection*{Mecânica principal}

A principal regra é:

\begin{center}
\textbf{NÃO SEJA VISTO.}
\end{center}

NPCs utilizam NavMesh para patrulhar corredores.

Cada inimigo possui:

\begin{itemize}
    \item rota de patrulha;
    \item campo de visão;
    \item estado de suspeita;
    \item estado de perseguição;
    \item retorno à patrulha.
\end{itemize}

\subsubsection*{Estados da IA}

\begin{center}
Patrulha
$\rightarrow$
Suspeita
$\rightarrow$
Alerta
$\rightarrow$
Perseguição.
\end{center}

Se o jogador escapar:

\begin{center}
Perseguição
$\rightarrow$
Busca
$\rightarrow$
Patrulha.
\end{center}

\subsubsection*{Descoberta narrativa}

Alex acessa um terminal.

O sistema apresenta arquivos sobre:

\begin{itemize}
    \item anatomia humana;
    \item atmosfera da Terra;
    \item doenças;
    \item cidades;
    \item tecnologia;
    \item armas;
    \item comportamento humano.
\end{itemize}

Alex percebe que os dinossauros estudam a Terra há muito tempo.

\subsubsection*{Meta de distribuição de assets -- planejamento}

Os números abaixo são metas de modelos visuais distintos para o
nível futuro, não um inventário de assets importados ou posicionados.
Repetições, mudanças de material e texturas isoladas não aumentam
essa contagem. A distribuição precisa ser comprovada na implementação.


\begin{center}
\begin{tabularx}{\textwidth}{|X|c|}
\hline
Arquitetura de laboratório & 18\\
\hline
Equipamentos científicos & 18\\
\hline
Computadores e terminais & 10\\
\hline
Câmaras e sistemas de contenção & 8\\
\hline
Objetos médicos & 7\\
\hline
Câmeras e segurança & 6\\
\hline
NPCs e elementos especiais & 6\\
\hline
\textbf{Total previsto} & \textbf{73}\\
\hline
\end{tabularx}
\end{center}

% =========================================================
% FASE 4
% =========================================================

\subsection{Nível 4 -- Os Rebeldes}

\textbf{Status: planejado, não implementado na versão atual.}


\subsubsection*{Sinopse}

Alex encontra uma pequena nave e tenta fugir.

O veículo é detectado e atacado.

Após a queda, Alex acorda em outra região de Kryos.

Um grupo de dinossauros o encontra.

Alex acredita inicialmente que será capturado novamente.

Porém, eles o ajudam.

São rebeldes contrários ao Projeto Retorno.

\subsubsection*{Objetivos}

\begin{enumerate}
    \item Sobreviver à queda.
    \item Explorar os destroços.
    \item Encontrar suprimentos.
    \item Seguir os rebeldes.
    \item Conhecer Tala.
    \item Descobrir a verdadeira situação de Kryos.
    \item Recuperar componentes.
    \item Ativar uma antiga estação.
    \item Preparar uma nova rota para a Terra.
\end{enumerate}

\subsubsection*{Grande revelação}

Tala mostra registros históricos.

Os ancestrais da população de Kryos viviam na Terra.

Quando detectaram o evento de extinção, uma pequena parte da
civilização conseguiu fugir.

Milhões de anos depois, seus descendentes estão tentando retornar ao
planeta de origem.

\subsubsection*{Conflito}

O governo acredita que possui direito histórico sobre a Terra.

Os rebeldes argumentam que a humanidade é agora a civilização
dominante do planeta e que um retorno militar criaria uma guerra
desnecessária.

\subsubsection*{Gameplay}

A fase combina:

\begin{itemize}
    \item exploração;
    \item diálogos;
    \item NPCs amigáveis;
    \item pequenos combates;
    \item recuperação de componentes;
    \item perseguição;
    \item utilização de NavMesh.
\end{itemize}

\subsubsection*{Meta de distribuição de assets -- planejamento}

Os números abaixo são metas de modelos visuais distintos para o
nível futuro, não um inventário de assets importados ou posicionados.
Repetições, mudanças de material e texturas isoladas não aumentam
essa contagem. A distribuição precisa ser comprovada na implementação.


\begin{center}
\begin{tabularx}{\textwidth}{|X|c|}
\hline
Natureza e vegetação & 20\\
\hline
Elementos de acampamento & 15\\
\hline
Destroços & 12\\
\hline
Equipamentos tecnológicos & 10\\
\hline
Ferramentas e equipamentos & 6\\
\hline
NPCs & 5\\
\hline
Objetos narrativos & 5\\
\hline
\textbf{Total previsto} & \textbf{73}\\
\hline
\end{tabularx}
\end{center}

% =========================================================
% FASE 5
% =========================================================

\subsection{Nível 5 -- Retorno à Terra}

\textbf{Status: planejado, não implementado na versão atual.}


\subsubsection*{Sinopse}

Alex consegue retornar à Terra.

Entretanto, naves de Kryos já começaram a aparecer.

A humanidade interpreta a aproximação como uma invasão.

Forças humanas e tropas de Rhaz entram em conflito.

O jogador precisa atravessar uma cidade parcialmente destruída para
chegar a um centro de comunicação.

\subsubsection*{Objetivos}

\begin{enumerate}
    \item Sair do local da queda.
    \item Atravessar a região urbana.
    \item Evitar ou enfrentar inimigos.
    \item Encontrar integrantes rebeldes.
    \item Recuperar dados sobre o Projeto Retorno.
    \item Chegar ao centro de transmissão.
    \item Enviar uma mensagem para ambos os lados.
    \item Confrontar Rhaz.
    \item Impedir o ataque principal.
\end{enumerate}

\subsubsection*{Gameplay}

É a fase mais orientada à ação.

O jogador utiliza:

\begin{itemize}
    \item cobertura;
    \item movimentação;
    \item arma simples;
    \item interação;
    \item NPCs;
    \item objetivos;
    \item checkpoints.
\end{itemize}

O combate continua sendo relativamente simples.

\subsubsection*{Objetivo final}

Alex precisa chegar a uma torre de comunicação.

Lá, utiliza os registros encontrados durante a campanha para transmitir
informações para humanos e dinossauros.

Os dados provam:

\begin{itemize}
    \item que Kryos realmente está morrendo;
    \item que existem grupos favoráveis à coexistência;
    \item que a invasão foi organizada apenas por uma parcela do governo;
    \item que uma negociação ainda é possível.
\end{itemize}

\subsubsection*{Final}

A frota interrompe temporariamente o ataque.

Rhaz é obrigado a negociar.

O jogo não mostra imediatamente a solução do conflito.

A cena final apresenta humanos e dinossauros iniciando o primeiro
contato diplomático oficial.

Alex observa enquanto representantes das duas espécies se aproximam.

A última frase é:

\begin{quote}
``Durante milhões de anos, acreditamos que pertencíamos a mundos
diferentes. Talvez estivéssemos apenas tentando encontrar o mesmo
lugar para chamar de lar.''
\end{quote}

A tela escurece.

\begin{center}
\textbf{FIM}
\end{center}

\subsubsection*{Meta de distribuição de assets -- planejamento}

Os números abaixo são metas de modelos visuais distintos para o
nível futuro, não um inventário de assets importados ou posicionados.
Repetições, mudanças de material e texturas isoladas não aumentam
essa contagem. A distribuição precisa ser comprovada na implementação.


\begin{center}
\begin{tabularx}{\textwidth}{|X|c|}
\hline
Prédios e estruturas urbanas & 20\\
\hline
Objetos urbanos & 15\\
\hline
Destroços & 10\\
\hline
Veículos & 8\\
\hline
Barricadas e equipamentos & 7\\
\hline
Personagens e NPCs & 6\\
\hline
Objetos tecnológicos & 7\\
\hline
\textbf{Total previsto} & \textbf{73}\\
\hline
\end{tabularx}
\end{center}

% =========================================================
% TUTORIAL
% =========================================================

\subsection{Níveis de treinamento}

Não há cena separada de tutorial. O desafio inicial da caixa apresenta
interação, mira, impulso e tentativa/erro dentro da própria contenção.

O HUD mostra o objetivo e as ações contextuais quando o jogador mira
em um objeto. Os menus oferecem o painel de controles. Não há, nesta
versão, uma sequência adicional de popups temporizados ensinando cada
tecla de movimentação.

O tutorial atual deve ser apresentado como parte do Nível 1, sem
declarar mensagens ou sistemas de ajuda ainda não implementados.



\section{Interface}

\subsection{Sistema visual}

A interface segue o sistema tradicional de UI da Unity, com
GameObject Canvas, Panel, Button, Text e eventos OnClick conectados
a scripts. Não utiliza UI Toolkit ou interface externa.

O Canvas usa \textit{Screen Space -- Overlay}; o Canvas Scaler usa
\textit{Scale With Screen Size}, com referência de 1920$\times$1080.
O estilo é futurista, com fundo escuro, detalhes azul/ciano e textos
claros. Há um EventSystem por cena de jogo/menu, usando o novo
Input System.

\subsubsection*{HUD de gameplay}

\begin{itemize}
\item Objetivo atual no canto superior esquerdo.
\item Mira pequena no centro da tela.
\item Ação contextual na parte inferior central.
\item Mensagens temporárias para coleta, instalação, adoção e ativações.
\item Progresso dos terminais, de 0/3 a 3/3.
\end{itemize}

Não há barra de vida, inventário visual ou indicador de detecção.
O HUD de gameplay é ocultado durante a pausa para não disputar
espaço com os menus.

\subsubsection*{Menu Principal}

A cena \texttt{MainMenu.unity} apresenta:

\begin{itemize}
\item Título: ECOS DO CRETÁCEO.
\item Subtítulo: PROJETO RETORNO.
\item Novo Jogo: inicia o Nível 1 com os estados zerados.
\item Controles: abre o painel dos comandos.
\item Sair: encerra a aplicação na build.
\end{itemize}

Não há botão de carregar save ou menu de opções na versão atual.

\subsubsection*{Menu de Pausa}

ESC abre ou fecha o painel sobre o cenário. O jogo usa
\texttt{Time.timeScale = 0}, mostra o cursor e bloqueia movimento,
câmera, interação e arremesso. O áudio de gameplay e as partículas
também pausam. A interface continua respondendo aos cliques.

Os botões são:

\begin{enumerate}
\item Continuar.
\item Controles.
\item Reiniciar Nível.
\item Menu Principal.
\item Sair.
\end{enumerate}

Reiniciar recarrega a cena atual. Trocar de cena restaura o tempo
normal. No Editor, Sair encerra o modo Play; na build, encerra
a aplicação.

\subsubsection*{Tela de conclusão}

Mostra o título NÍVEL 1 CONCLUÍDO, a confirmação da fuga e
o botão Voltar ao Menu Principal. Não oferece continuidade para
um Nível 2 ainda não implementado.

\subsection{Sistema de controle}

\begin{center}
\begin{tabularx}{\textwidth}{|p{4.3cm}|X|}
\hline
\textbf{Comando} & \textbf{Ação implementada}\\
\hline
W / A / S / D & Movimentação.\\
\hline
Mouse & Controlar câmera.\\
\hline
Espaço & Pular quando não estiver agachado.\\
\hline
Shift esquerdo & Correr.\\
\hline
Ctrl esquerdo ou direito, segurado & Agachar.\\
\hline
F & Interagir, coletar, usar painéis, pegar a primeira caixa e adotar o Sparrow.\\
\hline
Botão esquerdo do mouse & Arremessar uma caixa enquanto o desafio inicial estiver ativo.\\
\hline
ESC & Pausar / continuar.\\
\hline
\end{tabularx}
\end{center}

A tecla C não é um comando alternativo de agachamento.
O clique esquerdo não ataca ou dispara armas na versão atual.

\subsection{Áudio, música e efeitos sonoros}

Foram localizados 173 AudioClips dos pacotes Kenney Sci-fi Sounds
e Kenney Interface Sounds. O pacote adicional citado para falas,
anúncios e ambientes não foi encontrado importado.

Os clipes atualmente utilizados são:

\begin{center}
\begin{tabularx}{\textwidth}{|p{6.2cm}|X|}
\hline
\textbf{Clipe Kenney} & \textbf{Aplicação}\\
\hline
\texttt{Kenney\_doorOpen\_000.ogg} & Porta inicial e saída externa.\\
\hline
\texttt{Kenney\_doorOpen\_001.ogg} & Porta do laboratório.\\
\hline
\texttt{Kenney\_doorOpen\_002.ogg} & Porta do hangar.\\
\hline
\texttt{Kenney\_confirmation\_001.ogg} & Instalação bem-sucedida da bateria.\\
\hline
\texttt{Kenney\_error\_006.ogg} & Aviso eletrônico da saída externa, repetido por quatro segundos.\\
\hline
\texttt{Kenney\_engineCircular\_000.ogg} & Ambiente discreto da sala de energia.\\
\hline
\end{tabularx}
\end{center}

As fontes das portas e do pet são 3D. O som de abertura é disparado
quando a folha começa a mover, sem repetição por frame. O aviso
externo é finito, não uma fala ou sirene de um pacote inexistente.

Não há trilha musical implementada. A alternância entre músicas de
exploração, tensão e perseguição permanece como ideia para fases futuras.

\subsubsection*{Passos e cantos: sistemas preparados}

\texttt{PlayerFootstepAudio} lê velocidade real e contato com o chão.
A corrida aumenta a frequência; parado ou no ar não gera passos.
Pode variar pitch e escolher clipes sem repetição imediata.

\texttt{CompanionChirps} só agenda cantos após a adoção, em intervalos
aleatórios de 8 a 20 segundos, com volume discreto e fonte no pássaro.

\textbf{Os arrays estão vazios por falta de áudios adequados.}
Esses sons ainda não são ouvidos no jogo. Atribuir clipes reais no
Inspector completa a configuração, sem usar efeitos eletrônicos
como substitutos de passos ou cantos.

\subsection{Sistema de ajuda}

O objetivo orienta a progressão e a mensagem de interação descreve
a ação disponível. O painel Controles funciona no Menu Principal
e na pausa e possui botão Voltar.

Mensagens de requisito informam quando faltam bateria ou cartão.
As confirmações desaparecem após alguns segundos. Não há guia
por setas, marcador de missão no mundo ou diálogo de ajuda.



\section{Inteligência Artificial}

\subsection{Comportamento implementado do companheiro}

O Sparrow usa um comportamento simples e independente da progressão
dos puzzles. Antes da adoção, permanece no ponto de espera, com
variações periódicas de repouso. Depois de F, acompanha um alvo
lateral próximo ao jogador.

O movimento utiliza interpolação, verificação de obstáculos e
recuperação segura por distância ou bloqueio. Não existe patrulha,
combate, coleta de itens, navegação por NavMesh ou tomada de decisões
complexa no pet.

O pássaro pausa com o jogo, não bloqueia portas e não ativa os
triggers de conclusão. Seus clipes reais são controlados pelo
Animator original do pacote.

\subsection{IA oponente e inimiga -- planejamento}

A campanha prevê inimigos nos níveis posteriores, com máquina de
estados simples: repouso, patrulha, suspeita, alerta, perseguição,
busca e retorno.

O planejamento inclui pontos de patrulha, distância e campo de
visão, Raycast para verificar obstáculos e navegação por NavMesh.
Esses sistemas não estão implementados no Nível 1 e não devem ser
apresentados como uma funcionalidade atual das câmeras decorativas.

\subsection{NPCs amigáveis -- planejamento}

Tala e os rebeldes deverão participar dos níveis 4 e 5. Conversas,
acompanhamento narrativo, objetivos fornecidos por NPCs e navegação
por NavMesh são previstos para a campanha, não para o pet atual.

\subsection{Sistemas de suporte}

Na versão atual são utilizados Raycasts, Layers, colliders, triggers,
referências de componentes e interpolação. NavMesh e NavMesh Agent
continuam previstos para etapas futuras. O acompanhamento do
Sparrow não comprova a implementação desses dois recursos.



\section{Técnico}

\subsection{Hardware alvo}

A plataforma principal é Windows, com teclado e mouse. O objetivo
de apresentação é funcionar em computadores comuns do laboratório.

A referência de planejamento é Windows 10 ou superior, processador
Core i5 ou equivalente, 8 GB de RAM, GPU compatível com DirectX 11
e resolução a partir de 1280$\times$720. Esses valores não constituem
uma especificação mínima validada por testes de desempenho.

\subsection{Hardware e software de desenvolvimento}

\subsubsection*{Engine e renderização}

O projeto atual usa \textbf{Unity 6, versão 6000.0.81f1},
\textbf{Universal Render Pipeline (URP)}, C\# e o novo Input System.

O enunciado fornecido cita a versão \textbf{6000.0.8f1} do laboratório.
As versões não são idênticas. A compatibilidade e a versão aceita
para a apresentação precisam ser conferidas no laboratório; não
foi realizado downgrade ou alterada configuração global para esta
revisão do GDD.

\subsubsection*{Ferramentas}

São utilizados Unity Editor e código C\#, com assets importados e
ferramentas próprias de integração. Git/GitHub, Blender, editores
de áudio/imagem e outras ferramentas podem apoiar a equipe.
A existência de um asset em planejamento não comprova sua importação.

\subsubsection*{Cenas e fluxo}

\begin{itemize}
\item \texttt{Assets/Scenes/MainMenu.unity}: entrada da aplicação.
\item \texttt{Assets/Scenes/Level01\_Nave\_Expanded.unity}: versão jogável.
\item \texttt{Assets/Scenes/Level01\_Nave.unity}: referência original.
\end{itemize}

A Scene List mantém MainMenu primeiro. Novo Jogo carrega a nave
expandida; Reiniciar recarrega a cena atual. Não há progresso
persistente entre essas recargas.

\subsubsection*{Arquitetura dos scripts implementados}

O código é separado em \texttt{Runtime} e \texttt{Editor}.
Os nomes abaixo correspondem aos arquivos reais, não a classes
genéricas apenas planejadas.

\small
\begin{longtable}{|p{6.3cm}|p{8.7cm}|}
\hline
\textbf{Script ou grupo} & \textbf{Responsabilidade}\\
\hline
\endfirsthead
\hline
\textbf{Script ou grupo} & \textbf{Responsabilidade}\\
\hline
\endhead
\texttt{NaveWalkthroughController} & Movimento, câmera, gravidade, salto, corrida e agachamento existentes.\\
\hline
\texttt{PlayerInteraction} e \texttt{IInteractable} & Raycast, mensagem contextual e execução da interação F.\\
\hline
\texttt{Level01GameManager} & Estado único de bateria, energia, cartão, terminais, objetivos e conclusão.\\
\hline
\texttt{Level01Interactable} & Configuração comum para coleta, painel, leitor e terminais.\\
\hline
\texttt{DoorController} & Folha móvel, collider sólido, Translate e evento de início de abertura.\\
\hline
\texttt{ObjectiveHUD} & Objetivos, mira, mensagens, contador e tela final.\\
\hline
\texttt{LevelCompleteTrigger} & Presença da cápsula do jogador na saída, condicionada pelo gerente.\\
\hline
\texttt{CardSpawnSystem} & Sorteio do único cartão entre pontos válidos do laboratório.\\
\hline
\texttt{LaboratoryEntryTrigger} & Atualiza o objetivo ao entrar no laboratório.\\
\hline
\texttt{IntroThrowingSystem} & HoldPoint, array de prefabs, spawn por função, cooldown e desafio inicial.\\
\hline
\texttt{ThrownObject} & Rigidbody, impulso, torque, colisão contínua e vida útil da caixa.\\
\hline
\texttt{ThrowablePickup} e \texttt{ImpactButton} & Primeira coleta e validação do impacto no alvo.\\
\hline
\texttt{PetFollower} & Espera, adoção por IInteractable e acompanhamento seguro do Sparrow.\\
\hline
\texttt{SparrowAnimationController} & Idles aleatórios e transição para o Fly existente.\\
\hline
\texttt{CompanionChirps} & Cantos ocasionais após adoção, preparado sem clipes.\\
\hline
\texttt{PlayerFootstepAudio} & Cadência por movimento/chão, preparado sem clipes.\\
\hline
\texttt{PauseAwareAudio} & Pausa e retomada da fonte local de gameplay.\\
\hline
\texttt{DoorPresentation} & Som 3D, vapor e alerta finito, usando a abertura existente.\\
\hline
\texttt{EnergyActivationVFX} & Energia ciano e confirmação após instalação real da bateria.\\
\hline
\texttt{MainMenuController}, \texttt{PauseMenuController}, \texttt{ControlsPanelController}, \texttt{MenuSceneActions}
& Botões, controles, pausa e navegação entre cenas.\\
\hline
\texttt{NaveWorldLabel} & Atualização da fonte nas placas 3D do cenário.\\
\hline
\end{longtable}
\normalsize

O gerador \texttt{Level01NaveGenerator} mantém partes para estrutura,
decoração e validação. As ferramentas \texttt{Level01GameplaySetup},
\texttt{Level01AdjustmentsSetup} e \texttt{Level01PolishSetup}
integram gameplay e apresentação de modo organizado.

A integração de polimento pode ser reaplicada sem reconstruir a
nave. O gerador chama essas integrações em futuras reconstruções,
mas gerar novamente o mapa ainda pode substituir edições manuais.
Os ajustes recentes preservaram o layout existente e suas referências.

\subsubsection*{Organização dos arquivos}

\begin{itemize}
\item \texttt{Assets/Scripts/Runtime}: código usado pelo jogo.
\item \texttt{Assets/Scripts/Editor}: montagem, auditorias e testes.
\item \texttt{Assets/Scenes}: cenas e backups.
\item \texttt{Assets/Prefabs/Gameplay/Level01}: bateria, cartão e caixas locais.
\item \texttt{Assets/Audio/SFX}: Footsteps, Doors, Interactions, Alarms e Bird.
\item \texttt{Assets/Audio/Ambient}, \texttt{Music} e \texttt{Libraries}: ambiente, espaço para música e bibliotecas Kenney.
\item \texttt{Assets/VFX}: Particles, Materials e Prefabs.
\item \texttt{Docs}: referências, guias, inventários, testes e prévias.
\end{itemize}

Áudios foram organizados por \texttt{AssetDatabase.MoveAsset},
preservando metadados e GUIDs. Prefabs visuais dos fornecedores
foram reutilizados com referências reais. Adaptações de materiais
ficam em cópias locais compatíveis com URP.

\subsubsection*{Documentação das funções}

As funções C\# devem explicar finalidade, argumentos, retorno e
processamento. O código da equipe é mantido por responsabilidade,
para permitir estudo e apresentação pelos três integrantes.

\subsection{Requisitos acadêmicos e situação da entrega parcial}

O enunciado exige cinco níveis, com pelo menos 70 assets de
categorias diferentes por nível, e determina que o primeiro
bimestre implemente somente o Nível 1. Este GDD mantém os demais
níveis como planejamento.

\paragraph{Tema escolhido para a primeira etapa: Tecnologia}

O Nível 1 evidencia tecnologia por pelo menos cinco elementos
identificáveis: célula de energia, painel de instalação, leitor de
cartão, terminais de controle, portas automatizadas e câmeras de
vigilância decorativas. A função de cada elemento é descrita,
sem transformar repetições em novos modelos para a contagem de assets.

\paragraph{Demonstração técnica atual}

\small
\begin{longtable}{|p{6.3cm}|p{8.7cm}|}
\hline
\textbf{Critério citado no enunciado} & \textbf{Evidência e limite}\\
\hline
\endfirsthead
\hline
\textbf{Critério} & \textbf{Evidência e limite}\\
\hline
\endhead
Movimento e colisão & CharacterController, gravidade, corrida, salto, agachamento e sólidos da nave.\\
\hline
Translate & Abertura da folha móvel das portas.\\
\hline
AddForce, rotação e torque & Lançamento de caixa com impulso e AddTorque; câmera e personagem também giram.\\
\hline
Vetor de prefabs e escolha aleatória & Array de dois prefabs de caixas, selecionado por Random.Range.\\
\hline
Spawn por função & Instantiate na função de reposição das caixas; limite de projéteis e remoção por tempo.\\
\hline
Randomização de posição & Cartão único sorteado entre três apoios válidos no laboratório.\\
\hline
UI com texto, botão e reinício & Canvas nativo, HUD, menus, controles e Reiniciar Nível.\\
\hline
Som & Portas, confirmação da bateria, alerta externo e ambiente. Passos/cantos ainda sem clipes.\\
\hline
Partículas e fumaça/vapor & VFX ciano e vapor nativo nas portas. Não há explosão ou pedrinhas; confirmar a abrangência dos exemplos exigidos.\\
\hline
Animação & Pet com Idle\_A/B/C e Fly reais. Isso não comprova os estados parado, andando, correndo, pulando e caindo do personagem principal.\\
\hline
Tag & MainCamera configurada na câmera. Gameplay identifica interativos por componentes e Layers; não há Tags específicas de cartões ou inimigos. Conferir a demonstração esperada de Tag.\\
\hline
70 assets no Nível 1 & Catálogo-base formal com 82 modelos distintos; conferir inventário da cena final.\\
\hline
Scripts documentados & Funções com finalidade, entradas, retorno e processamento; código organizado para apresentação do grupo.\\
\hline
\end{longtable}
\normalsize

\paragraph{Pontos a validar, sem declarar cumprimento automático}

\begin{itemize}
\item Compatibilidade da versão instalada com a versão do laboratório.
\item Uso permitido de assets externos, conforme orientação do professor.
\item Cobertura dos estados de animação pedidos para o personagem;
movimentar a câmera ou animar o pet não substitui essa comprovação.
\item Abrangência dos exemplos de partículas e da utilização de Tag.
\item Inventário atualizado dos modelos da cena entregue.
\item Participação de cada integrante, avaliada proporcionalmente
pelo professor; este documento não atribui atividades não registradas.
\end{itemize}

\subsection{Planejamento do segundo bimestre}

O enunciado prevê a campanha completa, com cinco níveis,
NPCs, NavMesh, melhores animações, partículas criadas pelo grupo,
iluminação e pós-processamento.

Também descreve terreno modelado pela equipe e texturizado com
pelo menos dez layers, com elementos como morros, vales, água,
ravinas, areia, cascalho, grama, vegetação, árvores e caminhos.
Esses itens são planejamento de fases externas, não conteúdo
atual da nave.

O texto apresenta uma divergência sobre o tema: nos critérios do
segundo bimestre menciona três elementos de cada eixo; na descrição
geral reaparece o mínimo de cinco. A referência de planejamento
é cinco elementos por eixo, sujeita à confirmação do professor.
Ecorenovação e Skills ainda precisam de uma distribuição explícita
nos níveis futuros; não se declara que o simples uso de vegetação
ou tecnologia já comprova esses eixos.

Para a entrega são previstos GDD detalhado, arquivos completos
do projeto e apresentação no laboratório. O vídeo de 1 a 3 minutos,
com link, é indicado no enunciado somente para o segundo bimestre.
Essas exigências documentais não foram executadas automaticamente
por esta revisão do GDD.

\subsection{Validação registrada}

Os relatórios existentes registram compilação e execução na Unity
6000.0.81f1, progressão completa aprovada com 849 movimentos físicos
da cápsula e regressão dos menus aprovada em 21/10/2026.

Foram verificados portas fechadas, tetos nos nove setores,
Raycasts, estados negativos dos puzzles, sorteio do cartão,
adoção, voo, áudio, partículas, pausa, reinício e conclusão.

A auditoria de polimento comparou 1.434 renderizadores do cenário
ao backup e confirmou preservação de posição, rotação e escala.
O reposicionamento do pet para esperar no corredor foi intencional.
Isso não é uma contagem de modelos distintos.

Os 173 AudioClips tiveram GUIDs preservados. Reaplicar a integração
não acrescentou colisores auxiliares. Nenhum MeshCollider ou
Rigidbody novo foi adicionado à decoração nessa revisão.

Esses resultados são de testes no Editor. Não substituem teste no
laboratório, revisão auditiva humana ou validação de uma build
standalone, ainda não gerada nessa etapa.

\subsection{Requerimentos de rede}

O jogo é single-player e não utiliza servidor, multiplayer,
contas ou sincronização de progresso pela internet.



\section{Arte do Jogo}

\subsection*{Direção artística atual e planejada}

O Nível 1 utiliza uma direção sci-fi estilizada, modular e industrial,
com materiais metálicos, painéis, iluminação fria e detalhes ciano.
Assets de pacotes diferentes são combinados buscando coerência
visual e proporções adequadas à navegação humana.

Modelos e materiais originais dos fornecedores são preservados.
Quando necessária, a compatibilidade com URP é obtida por cópias
locais de materiais. Não se declara um estilo semi-realista completo
ou a produção própria de todos os modelos.

Para a campanha, continuam previstas três identidades:
a tecnologia de Kryos, os ambientes naturais do planeta e a região
urbana da Terra. Vegetação, terreno, água, ruínas externas, cidades
e dinossauros equipados pertencem ao planejamento dos níveis futuros,
não à comprovação do cenário atual.

\subsection*{Arquitetura e composição da nave}

A nave mantém salas conectadas, pisos contínuos, portas amplas e
corredores legíveis. Laboratório, depósito, energia, segurança,
hangar, máquinas e manutenção têm conjuntos de objetos e funções
distintos. A decoração não deve obstruir as rotas obrigatórias.

Computadores, mesas, caixas, contêineres, tubulações, estruturas
industriais, luminárias e câmeras formam a base visual.
O hangar possui uma saída com luz branca, sem planeta explorável
ou cenário exterior completo nesta etapa.

\subsection*{Sparrow e animações reais}

O pet reutiliza o prefab Sparrow -- Quirky Series, da Omabuarts
Studio, seu Animator Controller \texttt{AC\_Sparrow}, suas malhas e
seu atlas original, com material local compatível com URP.

As animações utilizadas são:

\begin{itemize}
\item \texttt{Idle\_A}, \texttt{Idle\_B} e \texttt{Idle\_C}: repouso
no chão antes da adoção, alternado em intervalos de 3 a 7 segundos.
\item \texttt{Fly}: voo após a adoção, com transição suave e sem
interrupções por poses terrestres.
\end{itemize}

O pacote também contém estados como Walk, Run e Sit, que não
são apresentados como usados pelo comportamento atual.
Existe um único clipe Fly utilizado; não é inventada uma coleção
de variações de voo.

As animações do companheiro não são equivalentes à implementação
das animações exigidas para o protagonista. Estados de personagem
humano, dano, morte, conversa e combate permanecem futuros.

\IfFileExists{Figuras/sparrow_companheiro.png}{
\begin{figure}[H]
\centering
\includegraphics[width=.95\textwidth]{Figuras/sparrow_companheiro.png}
\caption{Sparrow adotado acompanhando o jogador no corredor. Captura real da versão atual.}
\end{figure}
}{}

\subsection*{Partículas implementadas}

São usados Particle Systems nativos e materiais
\texttt{Universal Render Pipeline/Particles/Unlit}, com máscara
suave local e sem plugins externos.

\begin{itemize}
\item \texttt{VFX\_EnergyActivation}: pequenas partículas ciano à
frente do painel durante a restauração da energia, por cerca de
2 a 3 segundos.
\item \texttt{VFX\_DoorSteam}: vapor branco/cinza nas laterais das
quatro portas importantes, com descarga curta de aproximadamente
1 a 2 segundos.
\end{itemize}

Os prefabs são reutilizáveis, não têm colisores nem loop infinito.
A emissão é limitada e pausa com o jogo. O vapor não cobre toda a
passagem e a energia não é uma explosão.

Fogo, chuva, destroços, pedrinhas e explosões ainda não foram
implementados como efeitos da fase atual.

\IfFileExists{Figuras/energia_ativacao.png}{
\begin{figure}[H]
\centering
\includegraphics[width=.95\textwidth]{Figuras/energia_ativacao.png}
\caption{Partículas ciano na instalação da bateria, com objetivo e confirmação no HUD.}
\end{figure}
}{}
\IfFileExists{Figuras/vapor_porta.png}{
\begin{figure}[H]
\centering
\includegraphics[width=.95\textwidth]{Figuras/vapor_porta.png}
\caption{Descarga breve de vapor nas laterais da porta do laboratório durante a abertura.}
\end{figure}
}{}

\subsection*{Iluminação}

A iluminação da nave é compatível com URP e prioriza visibilidade
do percurso, tons frios e destaques tecnológicos, sem adição
desnecessária de sombras e luzes dinâmicas.

A luz branca da saída é formada por superfície local URP e luz
suave, ativadas somente depois dos três terminais. Não depende
de instalar Bloom ou um pacote de efeitos.

\begin{center}
\begin{tabularx}{\textwidth}{|c|X|}
\hline
\textbf{Nível} & \textbf{Sensação e situação}\\
\hline
1 & Fria, artificial e tecnológica; implementada na nave.\\
\hline
2 & Natural, ampla e desconhecida; planejada.\\
\hline
3 & Clínica, controlada e ameaçadora; planejada.\\
\hline
4 & Natural e acolhedora com áreas destruídas; planejada.\\
\hline
5 & Urbana e marcada pelo conflito; planejada.\\
\hline
\end{tabularx}
\end{center}

\subsection*{Evolução visual prevista}

Terrenos, vegetação, água, elementos urbanos, personagens
dinossauros e efeitos adicionais deverão ser selecionados ou
produzidos para os próximos níveis. Sua disponibilidade e adequação
serão conferidas em Assets, sem declarar pacotes apenas adquiridos
como já importados.

Qualidade visual, legibilidade, desempenho e possibilidade de
edição manual orientam as próximas decisões de arte.



\section*{Conclusão}

\addcontentsline{toc}{section}{Conclusão}

\jogo{} mantém a proposta de aventura narrativa em cinco níveis,
com mistério e revelações graduais sobre Kryos, a Terra e a
civilização dos dinossauros.

A entrega implementada concentra-se no Nível 1: sair da contenção
por um desafio de arremesso, restaurar energia, encontrar o cartão,
liberar o hangar, ativar três terminais e atravessar a saída iluminada.
O Sparrow oferece uma interação opcional de companheirismo sem
aumentar a complexidade dos puzzles.

O jogo atual preserva movimentação e menus, possui HUD, colisões,
spawn, randomização, áudio disponível, partículas e animações reais
do pet. As limitações de clipes, animação do protagonista e
compatibilidade com o laboratório estão registradas.

Os níveis 2 a 5 e seus sistemas narrativos, ambientais e de IA
continuam como planejamento. O documento diferencia essa intenção
da implementação existente, permitindo apresentar o trabalho
acadêmico com evidências e escopo verificável.



\end{document}
